% Avsnitt 13.7, ur lectures/L13/README.md.

\section{Sammanfattning}
\label{c13:sec:review}

\needspace{6\baselineskip}
\subsection*{Det här ska du kunna nu}
\begin{itemize}
  \item Implementera en bitseriell CRC-15-motor som matchar CAN:s polynom.
  \item Förklara varför generering och kontroll är samma operation, och vad ett nollställt
    CRC-register vid slutet av en mottagen ram bevisar.
  \item Räkna några bitar genom rekursionen för hand och jämföra med simuleringen.
  \item Säga varför motorn behöver nollställas mellan ramar, och vilket fel som uppstår om den inte
    gör det.
\end{itemize}

\needspace{6\baselineskip}
\subsection*{Frågor att testa dig själv med}
\begin{enumerate}
  \item Varför behöver motorn inget läge för ”generera” respektive ”kontrollera”?
  \item Vad betyder det att CRC-registret är noll efter att en mottagen ram matats igenom?
  \item Vad händer om en ram avbryts och nästa ram börjar matas in utan \code{clear}?
  \item Varför matas en instoppad bit inte in i CRC-motorn?
  \item Vilken signal i \code{can_controller} avgör när \code{crc15} får uppdatera sig, och varför
    är det inte \code{bit_done}?
\end{enumerate}

\needspace{6\baselineskip}
\subsection*{Riktvärde}
Vid det här laget bör två av de fyra lövmodulerna vara klara. Vilka två spelar ingen roll.

\needspace{6\baselineskip}
\subsection*{Blicken framåt}
\Chapref{c14:ch} bygger sändningsskiftregistret: att serialisera ut på bussen, med samma
stoppbitsregel som i \chapref{c10:ch}, nu i VHDL. \Chapref{c14:ch} och \ref{c15:ch} bygger
tillsammans de moduler som avgör vilka bitar som når \code{crc15} och när.
